% ======================================================================
% METHODS SECTION — for chapters/3-Theory-and-methods.tex
% Introduces coordinate system, probe indexing (Fres, Fx, Fy, Fz,
% M_RDT, P_in, P_out, P_bend, P0..P5) used throughout the results
% chapter, and the spectral analysis methodology.
% ======================================================================

\section{Transient simulation setup and monitor points}
\label{sec:monitor-setup}

A transient simulation of the baseline \ac{RDT}--\ac{GV}--duct
configuration was performed in ANSYS~CFX at $n = 650$~rpm to assess
the unsteady loading at the nominal operating point and to provide
the time histories needed for the spectral analysis. The interface
between the rotating \ac{RDT} domain and the stationary \ac{GV}
domain was specified as a Transient Rotor--Stator (TRS) interface.
Other solver settings follow those established for the steady-state
screening campaign (Section~\ref{sec:cfd-setup}).

The time step was fixed at $\Delta t = T_{\mathrm{rev}}/360
= 2.564\times 10^{-4}$~s, corresponding to one degree of rotor
rotation per time step. The total simulated physical time was
$15\,T_{\mathrm{rev}} \approx 1.38$~s, allowing the initial
transient to decay over approximately ten rotor revolutions before
the last five revolutions are retained for Fourier analysis.

\subsection{Coordinate system and reference positions}
\label{sec:coord-system}

The coordinate system used throughout the spectral analysis has its
origin at the geometric centre of the \ac{RDT} hub, with the
positive $z$-axis pointing downstream towards the \ac{RPT}. The
\ac{RDT} blade volume occupies
$-0.185~\mathrm{m} \leq z \leq +0.185~\mathrm{m}$, so that the
\ac{RDT}--\ac{GV} interface plane is located at
\begin{equation}
    z_{\mathrm{ref}} \;=\; -0.185~\mathrm{m}. \label{eq:zref}
\end{equation}
The duct radius is $R = 0.25$~m, consistent with the baseline
geometry established in~\cite{leonsen2025}. Downstream of
$z_{\mathrm{ref}}$ there is a short unloaded transition before the
\ac{GV} leading edge, and the \ac{GV} domain extends to
$z = -0.685$~m, giving a total \ac{GV} domain length of 500~mm.

Axial monitor positions are hereafter reported as the distance
downstream of the \ac{RDT}--\ac{GV} interface,
\begin{equation}
    s \;=\; |z| - |z_{\mathrm{ref}}|,
    \qquad 0~\mathrm{mm} \leq s \leq 500~\mathrm{mm},
\end{equation}
which is more convenient for physical interpretation than absolute
$z$-coordinates because $s$ counts from a physically meaningful
plane (the rotor--stator interface) rather than from the \ac{RDT}
hub centre.

\subsection{Monitor-signal indexing}
\label{sec:monitor-index}

Fifteen monitor signals are exported from the CFX solver manager and
used in the spectral analysis. To keep figure axis labels and table
rows compact, each signal is assigned a short label in
Table~\ref{tab:monitor-index}. These short labels are used on all
figures in Chapter~\ref{ch:results} and in Appendix~\ref{app:individual_spectra};
their full meaning is to be read from Table~\ref{tab:monitor-index}.

\begin{table}[H]
\centering
\small
\caption{Indexing of monitor signals used in the FFT analysis. The
short label is used on all plots. Positions refer to the coordinate
system defined in Section~\ref{sec:coord-system} ($s$ measured from
the \ac{RDT}--\ac{GV} interface at $z_{\mathrm{ref}} = -0.185$~m).}
\label{tab:monitor-index}
\begin{tabular}{llll}
\hline
Label & Quantity & Location / averaging & CFX expression \\ \hline
\multicolumn{4}{l}{\textit{Group A — forces on a single guide vane (GV1)}} \\
$F_{\mathrm{res}}$ & Resultant force     & Integrated on GV1 surface & \verb|Fres GV1| \\
$F_x$              & Force $x$-component & Integrated on GV1 surface & \verb|Fx GV1|   \\
$F_y$              & Force $y$-component & Integrated on GV1 surface & \verb|Fy GV1|   \\
$F_z$              & Force $z$-component & Integrated on GV1 surface & \verb|Fz GV1|   \\[0.4em]
\multicolumn{4}{l}{\textit{Group B — torque on the rotor}} \\
$M_{\mathrm{RDT}}$ & Torque about rotation axis & Integrated over full RDT
                                                & \verb|torque_z()@Entire BLADE| \\[0.4em]
\multicolumn{4}{l}{\textit{Group C — plane-averaged pressures}} \\
$P_{\mathrm{in}}$  & Static pressure, mass-flow avg.
                   & Plane at $s = 0$~mm (RDT--GV interface)
                   & \verb|massFlowAve(p)@RDT_GV| \\
$P_{\mathrm{out}}$ & Static pressure, mass-flow avg.
                   & Plane at $s = 500$~mm (GV outlet)
                   & \verb|massFlowAve(p)@GV_DT| \\
$P_{\mathrm{out}}^{\,a}$ & Static pressure, area avg.
                   & Plane at $s = 500$~mm (reference check)
                   & \verb|areaAve(p)@GV_DT| \\[0.4em]
\multicolumn{4}{l}{\textit{Group D — far-downstream centreline probe}} \\
$P_{\mathrm{bend}}$ & Static pressure & Pipe centreline, $z = -2.185$~m
                                      (before 90$^\circ$ bend)
                   & Monitor Point 1 \\[0.4em]
\multicolumn{4}{l}{\textit{Group E — wall-pressure probes at $r/R = 0.744$ along the GV passage}} \\
$P_0$ & Static pressure & $s = 0$~mm,   $r/R = 0.744$ (near GV leading edge) & \verb|pressure 185 GV| \\
$P_1$ & Static pressure & $s = 100$~mm, $r/R = 0.744$ (inside GV passage)    & \verb|pressure 285 GV| \\
$P_2$ & Static pressure & $s = 200$~mm, $r/R = 0.744$ (inside GV passage)    & \verb|pressure 385 GV| \\
$P_3$ & Static pressure & $s = 300$~mm, $r/R = 0.744$ (inside GV passage)    & \verb|pressure 485 GV| \\
$P_4$ & Static pressure & $s = 400$~mm, $r/R = 0.744$ (near trailing edge)   & \verb|pressure 585 GV| \\
$P_5$ & Static pressure & $s = 500$~mm, $r/R = 0.744$ (at GV outlet)         & \verb|pressure 685 GV| \\ \hline
\end{tabular}
\end{table}

\paragraph{Groups of purpose.}
The signals fall into four functional groups:
\begin{itemize}
  \item \emph{Group A — Forces on a single guide vane.} Four signals
        ($F_{\mathrm{res}}$, $F_x$, $F_y$, $F_z$) record the
        integrated hydrodynamic force on the surface of one
        representative guide vane (GV1). These are the principal
        diagnostics for cyclic structural loading of the vane.
  \item \emph{Group B — Torque on the rotor.}
        $M_{\mathrm{RDT}}$ is the torque about the rotation axis
        integrated over the entire \ac{RDT} blade set and captures
        periodic torque ripple produced by rotor--stator
        interaction.
  \item \emph{Group C — Plane-averaged pressures.} Two mass-flow-averaged
        plane signals, $P_{\mathrm{in}}$ at the \ac{RDT}--\ac{GV}
        interface and $P_{\mathrm{out}}$ at the \ac{GV} outlet, are
        used to extract the axisymmetric $m=0$ component of the
        pulsation field. The area-averaged downstream plane
        $P_{\mathrm{out}}^{\,a}$ is retained only as a reference check.
        Comparison of $P_{\mathrm{in}}$ and $P_{\mathrm{out}}$
        amplitudes quantifies how much of the coherent $m=0$ content
        is attenuated across the \ac{GV} row.
  \item \emph{Group D — Centreline probe before the bend.}
        $P_{\mathrm{bend}}$ serves as a check of how the pulsation
        field has decayed by the time the flow reaches the 90$^\circ$
        bend region.
  \item \emph{Group E — Wall-pressure probes along the GV passage.}
        Six probes $P_0 \ldots P_5$ at the fixed radial position
        $r/R = 0.744$ (approximately the half-area-radius,
        $r_{1/2} = R/\sqrt{2}$ gives $r/R = 0.707$) along one
        azimuthal line sample the full pulsation field, including
        the spinning modes $|m|\geq 1$ that are washed out in the
        plane-averaged signals. In combination with $P_{\mathrm{out}}$
        they permit the ring-coherence diagnostic defined in
        Eq.~\eqref{eq:coherence}.
\end{itemize}

\paragraph{Signals excluded from the analysis.}
Two additional static-pressure probes specified in the CFX setup at
$s = 600$~mm and $s = 700$~mm lie outside the updated \ac{GV}
domain; the downstream duct geometry was modified after the monitor
file was created, and their position therefore no longer corresponds
to a physically meaningful reference inside the \ac{GV} passage.
They are excluded from the spectral analysis. The area-averaged
pressure on the upstream plane is also excluded because the
mass-flow-averaged form is the correct projection onto the $m=0$
mode for a non-uniform velocity distribution.

\section{Spectral analysis methodology}
\label{sec:spec}

The force, torque and pressure signals from the monitors described
in Section~\ref{sec:monitor-index} are analysed in the frequency
domain using the discrete Fourier transform. This section summarises
the sampling, window and amplitude-calibration relations and defines
the ring-coherence diagnostic used to classify the dominant spatial
mode of the pulsation field.

\subsection{Discrete Fourier transform and sampling}
\label{sec:spec-dft}

A monitor signal $x(t)$ is recorded at the transient solver time
step $\Delta t$, giving $N$ uniformly spaced samples
$x_n = x(n\Delta t)$, $n = 0,1,\dots,N-1$. The sampling frequency
and Nyquist frequency are
\begin{equation}
    f_s \;=\; \frac{1}{\Delta t},\qquad
    f_{\mathrm{Nyq}} \;=\; \frac{f_s}{2}. \label{eq:fs_fnyq}
\end{equation}
The discrete Fourier transform and frequency grid are
\begin{align}
    X_k &\;=\; \sum_{n=0}^{N-1} x_n\, e^{-i\,2\pi k n / N},
    \qquad k = 0,1,\dots,N-1, \label{eq:dft}\\
    f_k &\;=\; \frac{k}{N\Delta t} \;=\; k\,\Delta f,
    \qquad \Delta f = \frac{1}{N\,\Delta t} = \frac{f_s}{N}.
    \label{eq:dft_f}
\end{align}
To place the excitation frequencies $f_n$ and $f_{\mathrm{BPF}}$
exactly on Fourier bins and thereby minimise spectral leakage, the
analysis window is chosen as an integer number of rotor revolutions,
\begin{equation}
    T_w \;=\; N_{\mathrm{rev}}\,T_{\mathrm{rev}},
    \qquad T_{\mathrm{rev}} = 1/f_n, \label{eq:Tw}
\end{equation}
so that $N = 360\,N_{\mathrm{rev}}$ at one degree of rotor rotation
per solver step. With $N_{\mathrm{rev}}=5$ this gives
$\Delta f = f_n/5$, i.e.\ five Fourier bins per rotor-frequency harmonic and
35 bins per blade passing harmonic.

\subsection{Hanning window and leakage correction}
\label{sec:spec-window}

Residual non-periodicity at the boundaries of the analysis window
introduces spectral leakage from dominant peaks into neighbouring
bins. Following standard practice for pressure-pulsation analysis
in hydraulic turbomachinery~\cite{kerlef,bergan2018} the signal is
multiplied by a Hanning window prior to the Fourier transform:
\begin{equation}
    w_n \;=\; \tfrac{1}{2}\!\left[ 1 - \cos\!\left(\tfrac{2\pi n}{N-1}\right)\right],
    \qquad n = 0,1,\dots,N-1. \label{eq:hanning}
\end{equation}
The mean is removed before windowing ($x'_n = x_n - \bar{x}$,
$x^{w}_n = x'_n w_n$). The Hanning window reduces the effective
signal energy by a factor $\sum_n w_n / N = 1/2$, so an amplitude
correction is applied to recover the true peak amplitude. The
single-sided amplitude spectrum used throughout this thesis is
\begin{equation}
    A_k \;=\; \frac{2}{\displaystyle\sum_{n=0}^{N-1} w_n}\,
      \bigl|X^{w}_k\bigr|,
    \qquad k = 0,1,\dots,\tfrac{N}{2}, \label{eq:amplitude}
\end{equation}
where the factor $2$ accounts for folding of negative frequencies
onto the single-sided representation. $A_k$ has the same physical
units as $x_n$ and represents the peak amplitude of a pure sinusoid
at frequency $f_k$.

\subsection{Periodicity verification}
\label{sec:spec-periodicity}

The Fourier amplitudes $A_k$ are only meaningful if the underlying
signal is statistically stationary over the analysis window. Two
verifications are therefore performed:
\begin{enumerate}
  \item \emph{Phase-folded comparison.} The $N_{\mathrm{rev}}$
        revolutions contained in the window are plotted on a common
        $[0,\,360^\circ]$ phase axis. A signal dominated by harmonics
        of $f_n$ produces nearly overlapping curves; deviation is a
        direct indicator of content at frequencies below $f_n$ or of
        incomplete transient decay.
  \item \emph{Window robustness test.} The spectrum is computed for
        $N_{\mathrm{rev}} = 5$ and $N_{\mathrm{rev}} = 6$ and the
        two spectra are overlaid. Stability of the peak amplitudes
        under variation of the window length confirms that the
        reported amplitudes are not an artefact of the particular
        window.
\end{enumerate}

\subsection{Ring-coherence diagnostic}
\label{sec:spec-coherence}

The Tyler--Sofrin analysis in Section~\ref{sec:rsi-tyler}
distinguishes between the axisymmetric $m=0$ mode and the spinning
modes with $|m|\geq 1$. Only the $m=0$ component produces a net axial
force on the \ac{GV} ring and propagates as a plane wave, whereas
the spinning modes cancel in circumferential integration.

To classify the dominant spatial mode of the computed pulsation field
without a full modal decomposition, a scalar ring-coherence
diagnostic is defined. Let $A_{\mathrm{loc}}(f)$ denote the Fourier
amplitude of a local wall-pressure signal (one of $P_0 \ldots P_5$)
and $A_{\mathrm{ring}}(f)$ the amplitude of the plane-averaged signal
$P_{\mathrm{out}}$ on the downstream plane. Because mass-flow
averaging extracts the $m=0$ component of the field,
$A_{\mathrm{ring}}$ measures the axisymmetric plane-wave content
alone. The ring-coherence ratio is defined as
\begin{equation}
    \mathcal{R}_m(f) \;=\;
    \frac{A_{\mathrm{ring}}(f)}{A_{\mathrm{loc}}(f)},
    \qquad 0 \leq \mathcal{R}_m(f) \leq 1. \label{eq:coherence}
\end{equation}
Two limiting interpretations follow:
\begin{itemize}
  \item $\mathcal{R}_m(f_{\mathrm{BPF}}) \to 1$ indicates that the
        pulsation at $f_{\mathrm{BPF}}$ is dominated by the $m=0$
        breathing mode. For $Z_r = Z_s$ this is the expected
        signature of the classical rotor-stator concern.
  \item $\mathcal{R}_m(f_{\mathrm{BPF}}) \ll 1$ indicates that the
        pulsation is dominated by spinning modes $|m|\geq 1$ that
        cancel in circumferential integration. The net axial ring
        force is then small and the classical $Z_r = Z_s$ concern
        does not manifest strongly in practice.
\end{itemize}

\subsection{Characteristic excitation frequencies}
\label{sec:spec-peaks}

Table~\ref{tab:characteristic_freqs} lists the characteristic
excitation frequencies for the baseline configuration at
$n = 650$~rpm with $Z_r = Z_s = 7$. These are the reference
frequencies marked in every spectrum in Chapter~\ref{ch:results}.

\begin{table}[H]
\centering
\caption{Characteristic excitation frequencies for the baseline
configuration ($n = 650$~rpm, $Z_r = Z_s = 7$).}
\label{tab:characteristic_freqs}
\begin{tabular}{lccc}
\hline
 & Symbol & Formula & Value [Hz] \\ \hline
Rotor frequency           & $f_n$               & $n/60$         & 10.83  \\
Second rotor-frequency harmonic     & $2f_n$              & $2n/60$        & 21.67  \\
Third rotor-frequency harmonic      & $3f_n$              & $3n/60$        & 32.50  \\
Blade passing frequency   & $f_{\mathrm{BPF}}$  & $Z_r\,f_n$     & 75.83  \\
Second BPF harmonic       & $2f_{\mathrm{BPF}}$ & $2Z_r\,f_n$    & 151.67 \\
Third BPF harmonic        & $3f_{\mathrm{BPF}}$ & $3Z_r\,f_n$    & 227.50 \\ \hline
\end{tabular}
\end{table}

A characteristic signature of rotor--stator interaction would be
dominant peaks at $f_{\mathrm{BPF}}$ and its harmonics. Dominant
energy at $f_n$ and its harmonics, in contrast, is typically
associated with a once-per-revolution forcing such as a rotor
imbalance, blade-to-blade variation in loading, or azimuthally
non-uniform inflow into the \ac{RDT}~\cite{brekke,bergan2018}. Both
types of content are analysed in Chapter~\ref{ch:results}.

% ---- Additional .bib entries used here ----
% @mastersthesis{leonsen2025,
%   author={Leonsen, H.H.},
%   title={Parametric hydraulic design of guide vanes for a rim-driven
%          thruster in pumped-storage hydropower (project report)},
%   school={NTNU}, year={2025}}
% @mastersthesis{kerlef,
%   author={Kerlef, M.},
%   title={Experimental study of reversible pump-turbine under steady
%          state operating conditions},
%   school={NTNU}, year={2025}}
% @article{bergan2018,
%   author={Bergan, C.W. and Solemslie, B.W. and Dahlhaug, O.G. and
%           Storli, P.-T.},
%   title={Hydraulic performance measurements of a Francis turbine},
%   journal={Flow Measurement and Instrumentation},
%   volume={63}, year={2018}}
